\documentclass[11pt]{article}

\usepackage[margin=1in]{geometry}
\usepackage{amsmath,amssymb,mathtools}
\usepackage{booktabs}
\usepackage{longtable}
\usepackage{array}
\usepackage{enumitem}
\usepackage{hyperref}
\usepackage[T1]{fontenc}
\usepackage{lmodern}
\usepackage{microtype}

\hypersetup{
  colorlinks = true,
  linkcolor = blue,
  urlcolor = blue,
  citecolor = blue
}

\newcommand{\RR}{\mathbb{R}}
\newcommand{\norm}[1]{\left\lVert #1 \right\rVert}
\newcommand{\filepath}[1]{\begingroup\urlstyle{tt}\nolinkurl{#1}\endgroup}
\newcommand{\code}[1]{\begingroup\urlstyle{tt}\nolinkurl{#1}\endgroup}
\newcommand{\vect}[1]{\mathbf{#1}}

\title{Standalone Specification for a Landmark Geodesic Kamada--Kawai Solver\\
\large v1 Prototype and v2 Fully Faithful Multiscale Target}
\author{Codex design note for the \filepath{/Users/pgajer/current_projects/grip} repository}
\date{March 30, 2026}

\begin{document}
\maketitle

\begin{abstract}
This document turns the geodesic Kamada--Kawai idea into an implementation
specification for the GRIP codebase. The immediate goal is a level-4 Sierpinski
carpet prototype that is explicit enough to implement with no ambiguity. The
longer-term goal is a fully faithful weighted multiscale variant that revisits
GRIP's insertion and warm-start strategy so that it is consistent with a
weighted geodesic energy. The document therefore specifies two versions:
\code{v1}, which keeps GRIP's current MISF construction and insertion warm
start unchanged and introduces the new landmark geodesic KK functional only in
the level-refinement phase, and \code{v2}, which replaces the legacy
hop-based insertion logic with a weight-aware strategy built from weighted
shortest-path anchors, barycentric or least-squares placement, and landmark
geodesic refinement on the active set.
\end{abstract}

\section{Purpose}

The purpose of this document is to make the new algorithm explicit before
implementation. In particular, it records:
\begin{itemize}[leftmargin=1.4em]
  \item the mathematical functional we want to optimize;
  \item the exact sparse variant chosen for implementation;
  \item the exact optimizer chosen for the first prototype;
  \item how the solver plugs into GRIP's current multiscale pipeline;
  \item which legacy pieces are intentionally retained in \code{v1};
  \item which pieces are deferred to a more faithful \code{v2}.
\end{itemize}

The guiding question is:
\begin{quote}
Can a landmark geodesic KK refinement preserve the global symmetry and
low-RMSE behavior we observed near the pre-final GRIP state, while improving
local smoothness without paying the full cost of all-pairs classical KK?
\end{quote}

\section{Why This Algorithm Is Worth Testing}

The current GRIP pipeline uses a multiscale MISF construction followed by
coarse-level modified KK refinement and a finest-level endgame that currently
uses FR by default and optionally \code{kk\_repulse}. The focused carpet
experiments carried out earlier in this repository showed:
\begin{itemize}[leftmargin=1.4em]
  \item the pre-final multiscale state can be globally very good;
  \item a tiny amount of finest-level smoothing helps;
  \item longer final FR schedules increase smoothness but also increase
    Procrustes RMSE on the level-4 carpet;
  \item structural modifications of the final FR stage bend the tradeoff in the
    right direction, but do not recover the near-KK RMSE regime.
\end{itemize}

This motivates testing a different final or level-refinement objective rather
than continuing to search the same FR family.

The proposed geodesic KK functional is attractive because it measures
discrepancy between graph-theoretic target distances and \emph{path length in
the embedded graph}, rather than straight-line chord length in the ambient
space. That is especially natural for weighted graphs, where the intended graph
metric is already encoded by edge weights.

\section{Current GRIP Multiscale Behavior}

The current GRIP multiscale insertion and refinement strategy matters because
\code{v1} intentionally keeps part of it unchanged.

\subsection{Current insertion and warm start}

When a finer MISF level is opened, GRIP inserts newly activated vertices one at
a time. The current implementation:
\begin{enumerate}[leftmargin=1.8em]
  \item performs a BFS from the new vertex over the graph;
  \item collects the first three already-placed higher-level vertices as
    anchors;
  \item assigns an initial position using either barycenter placement or a
    2D circle heuristic;
  \item performs three tiny local KK-style refinement steps against those same
    anchors;
  \item after all new vertices are inserted, refines the whole active set at
    that level.
\end{enumerate}

In the current codebase this logic lives primarily in:
\begin{itemize}[leftmargin=1.4em]
  \item \filepath{/Users/pgajer/current_projects/grip/src/MishSupport.cpp}
    (\code{bfs\_me\_v4()} and \code{KK\_spring\_local()}),
  \item \filepath{/Users/pgajer/current_projects/grip/src/DrawGraph.cpp}
    (\code{initial\_position\_barycenter()} and
    \code{initial\_position\_circle()}),
  \item \filepath{/Users/pgajer/current_projects/grip/src/MishEngine.cpp}
    (per-level opening and active-set refinement).
\end{itemize}

\subsection{Important limitation of the current warm start}

The current insertion warm start is not fully consistent with the proposed
geodesic weighted model:
\begin{itemize}[leftmargin=1.4em]
  \item it is built from hop-based BFS neighborhoods;
  \item the circle placement heuristic interprets anchor distances as Euclidean
    radii in the ambient space;
  \item the micro-refinement is a legacy local KK step rather than a geodesic
    energy descent.
\end{itemize}

This mismatch is acceptable for a first prototype if it is called out
explicitly and kept outside the definition of the new functional.

\section{Target Functional}

\subsection{Weighted graph metric}

Let $G=(V,E,w)$ be a connected weighted graph with vertex set
$V=\{1,\dots,n\}$ and positive edge weights $w_e > 0$ for all $e \in E$.
For vertices $i,j \in V$, let $g_{ij}$ denote the weighted shortest-path
distance in the graph metric induced by $w$.

Let $X = (\vect{x}_1,\dots,\vect{x}_n)$ with $\vect{x}_i \in \RR^d$ denote an
drawing of the active graph.

\subsection{Deterministic shortest-path choice}

For every ordered source vertex $i$, \code{v1} chooses a deterministic
shortest-path tree rooted at $i$:
\begin{itemize}[leftmargin=1.4em]
  \item run weighted Dijkstra from $i$;
  \item when two predecessor choices give equal shortest-path distance, choose
    the predecessor with the smallest vertex index;
  \item reconstruct $\gamma_{ij}$ by following the predecessor chain from $j$
    back to $i$.
\end{itemize}

This makes path selection explicit and reproducible. It also means the
functional is defined with respect to fixed chosen shortest paths rather than
the set of all possible equal-length shortest paths.

\subsection{Embedded geodesic length}

For an edge $e=\{u,v\}\in E$, define its embedded Euclidean length by
\[
|X(e)| = \norm{\vect{x}_u - \vect{x}_v}.
\]

Given the chosen graph shortest path $\gamma_{ij}$ from $i$ to $j$, define the
embedded geodesic path length by
\[
h_{ij}(X) = \sum_{e \in \gamma_{ij}} |X(e)|.
\]

\subsection{Target lengths and stiffness}

Let
\[
\ell_{ij} = L_0 g_{ij},
\qquad
L_0 = \frac{L}{D},
\qquad
k_{ij} = \frac{K}{\max(g_{ij}, \varepsilon_g)^2},
\]
where $D=\max_{i,j} g_{ij}$ is the weighted graph diameter on the active set,
$L$ is a drawing scale constant, $K$ is a global stiffness constant, and
$\varepsilon_g > 0$ is a small guard against division by zero.

\subsection{Landmark-sparse geodesic energy}

The all-pairs geodesic energy is conceptually clean but too expensive for the
first implementation. We therefore define a sparse constraint set
$\mathcal{S}_\ell$ on each active MISF level $\ell$.

For every active vertex $i$, define:
\begin{itemize}[leftmargin=1.4em]
  \item a local set $N_q(i)$ of the $q$ nearest active vertices to $i$ in the
    weighted graph metric;
  \item a landmark set $L_M(i)$ of size $M$, chosen by deterministic
    farthest-point sampling in weighted graph distance on the active set.
\end{itemize}

The v1 sparse pair set is
\[
\mathcal{S}_\ell
=
\bigl\{(i,j): i<j,\ j\in N_q(i)\cup L_M(i)
\text{ or } i\in N_q(j)\cup L_M(j)\bigr\}.
\]

The v1 landmark geodesic KK energy on the active set is
\[
E_{\mathrm{LGKK}}(X)
=
\frac{1}{2}\sum_{(i,j)\in \mathcal{S}_\ell}
k_{ij}\bigl(h_{ij}(X)-\ell_{ij}\bigr)^2.
\]

\section{Alternative Design Choices and Explicit v1 Choices}

Several choices could have been made differently. They are recorded here so the
implementation choice is explicit.

\begin{table}[h]
\centering
\caption{Key design alternatives and the chosen v1 decision.}
\label{tab:v1-choices}
\begin{tabular}{@{}>{\raggedright\arraybackslash}p{0.23\textwidth}>{\raggedright\arraybackslash}p{0.32\textwidth}>{\raggedright\arraybackslash}p{0.35\textwidth}@{}}
\toprule
Design axis & Alternatives considered & Chosen v1 decision \\
\midrule
Sparse long-range structure &
Global landmark pool shared by all vertices; per-vertex landmark sets $L(i)$ &
Use per-vertex landmark sets $L_M(i)$, because they are more faithful to the intended landmark-sparse functional and still manageable on the first carpet experiment. \\
\addlinespace
Local constraint set &
Immediate graph neighbors only; $q$ nearest active vertices in graph metric &
Use $q$ nearest active vertices in weighted graph distance to better match GRIP's multiscale local-neighborhood logic. \\
\addlinespace
Optimizer &
Majorization; pure gradient descent; projected descent; warm-start polish on top of GRIP &
Use warm-started gradient descent with backtracking line search, applied as the active-set level solver. \\
\addlinespace
Shortest-path choice &
Arbitrary shortest path; random tie-breaking; deterministic canonical tie-breaking &
Use deterministic canonical tie-breaking by smallest predecessor index. \\
\addlinespace
Integration point &
Apply only as a final-stage polish; apply on every active MISF level &
Apply on each active MISF level in place of the legacy level-refinement solver during the prototype experiment. \\
\bottomrule
\end{tabular}
\end{table}

\section{Gradient and Numerical Stabilization}

Let $(i,j)\in\mathcal{S}_\ell$ and let
\[
E_{ij}(X)=\frac{1}{2}k_{ij}\bigl(h_{ij}(X)-\ell_{ij}\bigr)^2.
\]
Suppose $\gamma_{ij}$ is represented as an ordered path
\[
v_0=i,\ v_1,\ \dots,\ v_m=j.
\]
Define a stabilized edge length
\[
\phi_\varepsilon(\vect{a},\vect{b})
=
\sqrt{\norm{\vect{a}-\vect{b}}^2 + \varepsilon_x^2},
\]
with a small constant $\varepsilon_x > 0$.
Then the stabilized path length is
\[
\tilde{h}_{ij}(X)
=
\sum_{t=0}^{m-1}\phi_\varepsilon(\vect{x}_{v_t},\vect{x}_{v_{t+1}}).
\]

The working implementation is allowed to optimize $\tilde{h}_{ij}$ instead of
$h_{ij}$ directly to avoid singular derivatives when two vertices coincide.

For a vertex $p$ on the path $\gamma_{ij}$, the pairwise gradient is
\[
\nabla_{\vect{x}_p} E_{ij}(X)
=
k_{ij}\bigl(\tilde{h}_{ij}(X)-\ell_{ij}\bigr)
\sum_{e=\{u,v\}\in\gamma_{ij}:\,p\in\{u,v\}}
\sigma_e(p)\frac{\vect{x}_u-\vect{x}_v}
{\phi_\varepsilon(\vect{x}_u,\vect{x}_v)},
\]
where
\[
\sigma_e(p)=
\begin{cases}
1, & p=u,\\
-1, & p=v.
\end{cases}
\]

If $p$ is not on $\gamma_{ij}$, then $\nabla_{\vect{x}_p}E_{ij}(X)=0$.

The full gradient is
\[
\nabla E_{\mathrm{LGKK}}(X)
=
\sum_{(i,j)\in\mathcal{S}_\ell} \nabla E_{ij}(X).
\]

\section{Optimizer}

\subsection{Chosen v1 optimizer}

The v1 optimizer is a warm-started deterministic gradient descent with
backtracking line search:
\begin{enumerate}[leftmargin=1.8em]
  \item start from the current active-set coordinates inherited from the GRIP
    multiscale process;
  \item compute the full sparse gradient of $E_{\mathrm{LGKK}}$;
  \item if the gradient norm is below tolerance, stop;
  \item propose a step along $-\nabla E$;
  \item use backtracking line search until the energy decreases;
  \item accept the step, recenter the active layout to zero mean, and continue.
\end{enumerate}

The implementation will use fixed numerical defaults for:
\begin{itemize}[leftmargin=1.4em]
  \item maximum iterations per level;
  \item Armijo decrease factor;
  \item line-search shrink factor;
  \item gradient-norm tolerance;
  \item stabilization constants $\varepsilon_g$ and $\varepsilon_x$.
\end{itemize}

\subsection{Why this optimizer was chosen}

This choice is deliberate:
\begin{itemize}[leftmargin=1.4em]
  \item it is straightforward to implement from the explicit gradient;
  \item it is deterministic;
  \item it can be warm-started from existing GRIP coordinates;
  \item it isolates the new energy from unrelated changes in the insertion
    logic;
  \item it is good enough for the first level-4 carpet experiment.
\end{itemize}

Majorization or more specialized stress solvers may still be better long-term,
but they are intentionally deferred.

\section{Multiscale Integration}

\subsection{v1 integration}

The v1 prototype is intentionally conservative.

\paragraph{Explicit v1 statement.}
\emph{Version 1 retains GRIP's existing MISF construction and insertion warm
start unchanged.} In particular, \code{v1} does \textbf{not} alter:
\begin{itemize}[leftmargin=1.4em]
  \item MISF construction;
  \item BFS-based anchor discovery during insertion;
  \item barycenter placement;
  \item circle-heuristic placement;
  \item the legacy local KK micro-polish used immediately after insertion.
\end{itemize}

In \code{v1}, the new landmark geodesic KK solver enters only as the
\emph{active-set level-refinement step}. This choice is made to isolate the new
functional in the first experiment.

\subsection{v2 target}

\paragraph{Explicit v2 statement.}
\emph{Version 2 is intended to be a fully faithful weighted multiscale
implementation.} It will revisit insertion and replace the legacy local KK warm
start with a weight-aware alternative, likely built from:
\begin{itemize}[leftmargin=1.4em]
  \item weighted shortest-path anchors;
  \item barycentric or least-squares anchor placement;
  \item no circle heuristic;
  \item landmark geodesic KK refinement on the active set.
\end{itemize}

The purpose of this v2 statement is to make it explicit that the v1 insertion
pipeline is an engineering scaffold, not the intended final multiscale design.

\paragraph{Asymmetry propagation note.}
The level-4 carpet contact sheets generated during the current investigation
suggest a specific failure mode worth tracking into later development: the best
current GRIP-like layouts can retain a visibly skewed central void, whereas the
\code{igraph::layout\_with\_kk()} reference produces a much more symmetric inner
square. A plausible explanation is that asymmetry introduced during multiscale
insertion propagates forward and is only partially corrected by later
refinement. For that reason, if \code{v2} still shows central-hole skew after
the geodesic KK solver is in place, insertion symmetry should be treated as a
first-class debugging target rather than as a minor polishing issue.

\section{Algorithm Specification}

\subsection{v1 prototype algorithm}

For each MISF level from coarse to fine:
\begin{enumerate}[leftmargin=1.8em]
  \item activate the next MISF level using the current GRIP machinery;
  \item insert new vertices using GRIP's existing insertion routine unchanged;
  \item form the active-set weighted graph;
  \item compute weighted shortest-path distances on the active set;
  \item build local sets $N_q(i)$ and landmark sets $L_M(i)$ for active
    vertices;
  \item assemble the sparse pair set $\mathcal{S}_\ell$ and the chosen shortest
    paths $\gamma_{ij}$;
  \item initialize the optimizer at the current active coordinates;
  \item minimize $E_{\mathrm{LGKK}}$ on the active set by deterministic
    gradient descent with backtracking;
  \item continue to the next finer level.
\end{enumerate}

\subsection{Initial parameter defaults for the first prototype}

The first prototype should expose these explicit defaults:
\begin{itemize}[leftmargin=1.4em]
  \item dimension: $d=2$ for the level-4 carpet experiment;
  \item local-neighborhood size: $q = 20$;
  \item per-vertex landmark count: $M = 8$ or $M = 16$ in the first sweep;
  \item optimizer iterations per level: initially modest, e.g.\ 8, 16, 32;
  \item deterministic tie-breaking everywhere.
\end{itemize}

These defaults are not claimed to be optimal. They are chosen to make the first
experiment concrete and bounded.

\section{Evaluation Protocol}

Before the optimizer itself is trusted, the implementation should include an
energy evaluator for $E_{\mathrm{LGKK}}$ and use it to score known layouts on
the level-4 carpet:
\begin{itemize}[leftmargin=1.4em]
  \item canonical embedding;
  \item true pre-final GRIP state;
  \item one-round final FR state;
  \item best Round-3 structural GRIP layout;
  \item \code{igraph::layout\_with\_kk()}.
\end{itemize}

This is a ranking sanity check. If the new energy does not rank these layouts
in a sensible way, the optimization stage should be revisited before deeper
engineering.

After the evaluator check, the first full optimization experiment should be:
\begin{itemize}[leftmargin=1.4em]
  \item graph family: Sierpinski carpet;
  \item graph size: level 4;
  \item comparison baselines: current GRIP, current best structural GRIP,
    \code{igraph::layout\_with\_kk()};
  \item metrics: Procrustes RMSE, edge-length CV, axis deviation, boundary
    waviness, corridor waviness, hole-center error, runtime.
\end{itemize}

Immediately after the carpet test, at least one genuinely weighted graph test
should be added, because the weighted-graph motivation is a central reason for
this algorithm.

\section{Success Criteria}

The first prototype is successful if it satisfies all of the following:
\begin{itemize}[leftmargin=1.4em]
  \item the energy evaluator is deterministic and numerically stable;
  \item the optimizer decreases $E_{\mathrm{LGKK}}$ reliably on the level-4
    carpet;
  \item the resulting layouts are competitive with current GRIP on the carpet;
  \item the runtime remains far below a full all-pairs classical KK solve on
    the same graph;
  \item the code path is explicit enough that v2 can later replace the legacy
    insertion logic without redesigning the energy.
\end{itemize}

\section{Open Items Explicitly Deferred}

The following are intentionally out of scope for v1:
\begin{itemize}[leftmargin=1.4em]
  \item replacing MISF construction itself with a weight-aware filtration;
  \item revisiting insertion anchors and removing the circle heuristic;
  \item using majorization or second-order optimization;
  \item averaging over multiple equal shortest paths instead of choosing one
    deterministic path;
  \item extending the first prototype beyond the bounded carpet-first
    experiment.
\end{itemize}

\section{Recommendation}

This document recommends the following implementation order:
\begin{enumerate}[leftmargin=1.8em]
  \item implement the sparse geodesic energy evaluator first;
  \item validate its ranking behavior on the level-4 carpet;
  \item implement the v1 level-refinement optimizer with current GRIP insertion
    retained unchanged;
  \item benchmark it against current GRIP and \code{igraph::layout\_with\_kk()};
  \item only then begin the v2 work on weighted-consistent insertion.
\end{enumerate}

That order gives a clean, controlled first experiment while keeping the
longer-term fully faithful weighted multiscale target explicit from the start.

\end{document}
